% Generated draft: requires mathematical and source review.
\documentclass[11pt]{article}
\usepackage[margin=25mm]{geometry}
\usepackage{fontspec}
\setmainfont{Latin Modern Roman}
\usepackage{amsmath,amssymb,mathtools,unicode-math}
\setmathfont{Latin Modern Math}
\usepackage{enumitem,needspace,xurl,hyperref}
\setlength{\parindent}{0pt}
\setlength{\parskip}{5pt}
\setlength{\emergencystretch}{3em}
\newcommand{\N}{\mathbb{N}}
\newcommand{\Z}{\mathbb{Z}}
\newcommand{\Q}{\mathbb{Q}}
\newcommand{\R}{\mathbb{R}}
\newcommand{\C}{\mathbb{C}}
\newcommand{\F}{\mathbb{F}}
\newcommand{\A}{\mathbb{A}}
\newcommand{\PP}{\mathbb{P}}
\newcommand{\E}{\mathbb{E}}
\newcommand{\eps}{\varepsilon}
\newcommand{\dd}{\,\mathrm{d}}
\begin{document}
\section*{UnsolvedMath ID: 20000321}
\textbf{Source problem code:} catalogue entry
\subsection{Definitions and mathematical statement}
A stable surface is conventionally a projective surface with semi-log-canonical singularities and ample $\Q$-Cartier canonical divisor. Fix the relevant bounded moduli data, including the source's $K^2$ and $\chi(\mathcal O)$; the precise index/Hilbert data are not supplied. A smoothable surface occurs as a special fiber of an appropriate family with smooth or rational-double-point general fibers, but the source asks for a more precise infinitesimal notion rather than choosing one here.

The recovered 2004 Problem 1.2 has two parts. First, give a satisfactory boundary scheme structure on $\overline{M^{\mathrm{sm}}_{K^2,\chi}}$, the compactification of the smoothable locus. The source says discarding components parametrizing only worse-than-rational-double-point surfaces loses infinitesimal information. Second, retaining all stable-surface components, ask whether every punctured-disk family $\mathcal X\to\Delta'$ satisfying Koll\'ar's condition extends, after possible base change, to a stable family over $\Delta$ satisfying the same condition. Koll\'ar's condition concerns compatible relative reflexive canonical powers and base change; its exact flatness/index formulation is not included. In the usual notation $\omega^{[m]}=(\omega^{\otimes m})^{**}$ for positive $m$.

The supplied report describes a characteristic-zero KSBA answer through the scheme-theoretic closure of a rational-Gorenstein open substack $U$ in the full stack $\mathcal M$, after fixing precise bounded data. It attributes stable reduction after finite base change to Hacon--Xu, and uniqueness up to unique isomorphism under a finite separated diagonal hypothesis. The exact interior, relation to the original smoothability notion, and passage from stack closure to a coarse boundary scheme require specification.

Its candidate formal lemma is: for a noetherian algebraic stack $\mathcal M$ and open immersion $j:U\hookrightarrow\mathcal M$, let $\mathcal Z$ be cut out by
\[
 \ker(\mathcal O_{\mathcal M}\longrightarrow j_*\mathcal O_U).
\]
It claims $\mathcal Z$ is the smallest closed substack containing $U$ schematically densely, inherits properness over a base \emph{if} $\mathcal M$ is proper over that base, and every map from an integral discrete valuation ring to $\mathcal M$ whose generic point maps to $U$ factors uniquely through $\mathcal Z$. This would place a full-stack stable extension with such generic fiber in the scheme-theoretic smoothable closure. It does not make an arbitrary noetherian stack proper or identify the closure with a reduced union of components.

For positive/mixed characteristic, the report mentions fixed-volume stacks over $\Z[1/30]$ and sufficiently-large-characteristic stable reduction for fixed data, while explicitly leaving uniform properness over $\Z[1/30]$ conditional on locally stable reduction and calling small characteristics delicate. These later assessments and the formal lemma remain unreviewed reports, distinct from the original source-era lack-of-proof statements.
\subsection{Short English statement}
Give an infinitesimally meaningful smoothable boundary and stable reduction for all components while preserving Kollár's condition. The report proposes characteristic-zero KSBA closure/reduction results and a general DVR factorization lemma. It keeps the positive/mixed-characteristic properness claims conditional.
\subsection{Sources}
\begin{itemize}
\item Local UnsolvedMath registry, permanent ID $20000321$, source code \texttt{catalogue entry}; created and updated 2026-08-14. Status \texttt{solved}; classification \texttt{SOLVED-IN-LITERATURE}; difficulty L3; no proposer or year supplied.
\item Michael A. van Opstall's notes, AIM workshop \emph{Compact moduli spaces and birational geometry}, Palo Alto, December 6--10, 2004, Problem 1.2, source type PDF, \url{https://aimath.org/WWN/birational/birational.pdf}. Canonical location \texttt{aim-algebraic-geometry-notes.json}, \texttt{notes[320]}, tag \texttt{problem}; statement status \texttt{reconstructed\_unverified}; clean statement null. Verification attributes recovery of $\overline{M^{\mathrm{sm}}_{K^2,\chi}}$ and $\Delta'$ to official source TeX, and identifies ``12MICHAEL A. VAN OPSTALL'' as a damaged page header. Original remarks and literature field empty.
\item The original says Alexeev alluded to the issues in a talk. The research report invokes modern KSBA theory and Hacon--Xu stable reduction, without exact primary bibliography or bounded-data hypotheses.
\item Supplied research attempt 2: result \texttt{partial\_result}, run status \texttt{solved}, rigor label \texttt{proved}, low novelty confidence. Candidate contribution \texttt{lemma}: scheme-theoretic closure and integral-DVR factorization. Explicit caveat: machine-generated research aid, not peer reviewed; independently verify before citation.
\item Dataset attribution: Ulam AI / UnsolvedMath Contributors, CC BY 4.0; original problem attribution remains to the supplied workshop notes.
\end{itemize}
\end{document}
